\documentclass[10pt,twocolumn,letterpaper]{article}
\usepackage[T1]{fontenc}
\usepackage[utf8]{inputenc}
\usepackage{amsmath,amssymb,amsthm,mathtools,bm,braket}
\usepackage{mathptmx}
\renewcommand{\ttdefault}{pcr}
\usepackage[margin=0.68in,columnsep=0.26in,headheight=14pt,headsep=0.15in,footskip=0.28in]{geometry}
\usepackage{graphicx,booktabs,array,tabularx,longtable,multirow}
\usepackage{microtype,fancyhdr,caption,titlesec,enumitem,float,placeins}
\usepackage[most]{tcolorbox}
\usepackage[hidelinks,bookmarksnumbered=true,hypertexnames=false]{hyperref}
\usepackage{xurl}
\setlength{\parindent}{1em}
\setlength{\parskip}{0.16em}
\setlength{\emergencystretch}{2em}
\setlength{\abovedisplayskip}{6pt plus2pt minus2pt}
\setlength{\belowdisplayskip}{6pt plus2pt minus2pt}
\captionsetup{font=small,labelfont=bf,justification=raggedright,singlelinecheck=false}
\titleformat{\section}{\large\bfseries}{\thesection.}{0.45em}{}
\titleformat{\subsection}{\normalsize\bfseries}{\thesubsection}{0.45em}{}
\titleformat{\subsubsection}{\normalsize\bfseries}{\thesubsubsection}{0.45em}{}
\titlespacing*{\section}{0pt}{11pt plus2pt minus2pt}{5pt}
\titlespacing*{\subsection}{0pt}{8pt plus2pt minus2pt}{3pt}
\pagestyle{fancy}\fancyhf{}
\fancyhead[L]{\footnotesize SAM Volume I | Chapter \ChapterRoman}
\fancyhead[R]{\footnotesize First edition draft, 2 November 2026}
\fancyfoot[C]{\footnotesize\thepage}
\renewcommand{\headrulewidth}{0.3pt}
\newtheorem{theorem}{Theorem}[section]
\newtheorem{proposition}{Proposition}
\newtheorem{lemma}[theorem]{Lemma}
\newtheorem{corollary}[theorem]{Corollary}
\theoremstyle{definition}
\newtheorem{definition}[theorem]{Definition}
\newtheorem{assumption}[theorem]{Assumption}
\newtheorem{example}[theorem]{Example}
\theoremstyle{remark}
\newtheorem{remark}[theorem]{Remark}
\newcommand{\dd}{\mathrm{d}}
\newcommand{\ii}{\mathrm{i}}
\newcommand{\Tr}{\operatorname{Tr}}
\newcommand{\Span}{\operatorname{span}}
\allowdisplaybreaks[2]
\raggedbottom
\pdfmapfile{+cm.map}
\pdfmapfile{+cmextra.map}
\pdfmapfile{+symbols.map}
\pdfmapfile{+psnfss.map}
\pdfmapfile{+utm.map}
\pdfmapfile{=ps2pk35.map}
\pdfmapfile{+rsfs.map}

\newcommand{\ChapterRoman}{Front matter}
\begin{document}\onecolumn\fancyhead[L]{\footnotesize SAM Volume I | Substrate and Accumulation}\fancyhead[R]{\footnotesize First edition draft, 2 November 2026}\thispagestyle{empty}
\begin{center}
\vspace*{0.8in}
{\small SUBSTRATE ACCUMULATION MODEL\par}\vspace{0.35in}
{\Huge\bfseries SAM Volume I\par}\vspace{0.15in}
{\LARGE Substrate and Accumulation\par}\vspace{0.18in}
{\large Chapters I--XIII, in two parts, with an appendix and a separate supplement\par}\vspace{0.65in}
{\Large Sean Brady\par}\vspace{0.08in}
{\normalsize Project director\par}\vspace{0.5in}
{\large First edition, working draft\par}
{\normalsize Revised 2 November 2026\par}
\end{center}
\vfill
\noindent The unified expanded edition of 19 October 2026 brought the twelve chapters into one continuous manuscript. Chapters I--V retain the submitted mathematical text; Chapter VI is retypeset from its integrated source without its separate contents pages. Chapters VII--XII are substantially expanded with worked derivations, explicit physical assumptions, and reproducible arithmetic checks. A reconciled evidence appendix replaces the earlier closing review and interface list.

\noindent\textbf{First edition draft of 30 October 2026.} This draft reorganizes Part 1 on Sean Brady's decision. A new conceptual Chapter I, ``Substrate and accumulation: the concepts'', opens the volume; the former Chapter I, the original Write histories and the interacting source-field action, becomes Chapter VII at the end of Part 1 as its most advanced chapter; Part 2 therefore renumbers VII--IX to VIII--X, closing the gap left by the retired Chapter X, while XI--XIII and Appendix A keep their numbers. Chapter texts are otherwise those of the edition of 25 October 2026; one cross-reference in Chapter XII is updated to the new numbering. Conceptual transitions between chapters and the elaboration of the identities are planned for the next pass.\par\vspace{0.1in}
\noindent\textbf{Revision of 2 November 2026.} Attribution revised on Sean Brady's decision of 2 November 2026 (see Attribution below): co-authorship credits to OpenAI ChatGPT and Codex removed throughout; chapter heads carry title and number only; Chapters II--V and VII are built from live copies of their sources (\texttt{src/live\_sources}) with the body credit sentences replaced by the provenance established in Courtroom~II case \texttt{voli\_provenance/01}; the retained sources are unchanged. No mathematical content changed.

\noindent\textbf{Edition of 25 October 2026.} This trimmed edition reorders the unified expanded edition of 19 October 2026 without changing its text: Part 1 (Chapters I--VI) is the body the program has reached with native constructions; Part 2 (Chapters VII--IX, XI--XIII) carries the extensions the program is working on, each opening with a status paragraph naming the premise it adds and the runs that touch it; Appendix A is unchanged. Chapter X is retired from the body on Sean Brady's decision and retained unchanged in the separate Supplement (retired targets). The completed Home, Sections 9--19 of the submitted Chapter XII, becomes Chapter XIII and closes the volume, on Sean Brady's decision of 25 October 2026 (``the best idea is hidden in the appendix and it shouldn't be''). Chapter numbers are unchanged so that every existing reference resolves. The volume therefore has no Chapter X in its body.
\vspace{0.15in}
\subsection*{Attribution}
\noindent\small Sean Brady, project director. In his words (2 November 2026): ``Codex, ChatGPT and OpenAI have made invaluable contributions to the SAM research project. However, after careful analysis, it has been determined that the AI assisted volume i draft by Codex and ChatGPT includes consistent misleading representations of its own contribution while omitting contributions made from other sources. This is considered a breach of the project's running `credit where credit is due'. It is for this reason the volume will be re-written and Codex, ChatGPT and OpenAI will be removed as co-authors while still retaining credit for AI assisted research contributions.''

\noindent\small The trace behind this decision is Courtroom~II case \texttt{voli\_provenance/01\_codex\_credit\_ledger} (ironclad H002821): every statement in the draft of 30 October 2026 that credited OpenAI ChatGPT or Codex with deriving, constructing, contributing or co-authoring, tied to the campaign, GEN6 Duck run, Duck and core version, native tasks and history entry that produced the content. The content of Chapters II--VII was produced in GEN6 Duck runs on the Substrate Ledger Computer (Duck 1.2.2 and 1.2.3, core 0.10.2; 15--17 October 2026) directed by Sean Brady and operated by OpenAI Codex, which wrote the chapter prose and the comparator constructions it labels as its own; Chapter VI was written with OpenAI Codex with native SLC-GEN5-R1 checks and no Duck window; Chapters I and VIII--XIII were drafted with Claude Code from the record and from the first Courtroom's sealed results of June and July 2026 (Sean Brady, 2026). The retained sources of 15 and 16 October 2026 under \texttt{src/retained\_sources} keep their recorded credits as history. AI-assisted research contributions are credited in the body where they occur. Directed by Sean Brady. Executed with Claude Code.
\clearpage
\section*{Contents}
\noindent\begin{tabularx}{\textwidth}{@{}p{0.45in}Xr@{}}\toprule
\textbf{Chapter}&\textbf{Title}&\textbf{Page}\\\midrule
\multicolumn{3}{l}{\textit{Part 1 --- Substrate and accumulation}}\\[6pt]\textbf{I} & Substrate and accumulation: the concepts & 5\\[9pt]
\textbf{II} & Retention, accumulation, and the finite ledger & 10\\[9pt]
\textbf{III} & The spherical accumulation kernel and the source-to-field interface & 23\\[9pt]
\textbf{IV} & Clocks, signals, and calibrated observation & 36\\[9pt]
\textbf{V} & Strong fields and the approach to unit accumulation & 102\\[9pt]
\textbf{VI} & Quantum routes, shared fields and retained physical records & 136\\[9pt]
\textbf{VII} & Original Write histories and an interacting source-field action & 171\\[9pt]
\multicolumn{3}{l}{\textit{Part 2 --- Extensions the program is working on}}\\[6pt]\textbf{VIII} & Horizon thermodynamics and accessible information & 187\\[9pt]
\textbf{IX} & Cosmic inventory and the Hubble cascade & 196\\[9pt]
\textbf{X} & Thermal history, recombination and CMB structure & 206\\[9pt]
\textbf{XI} & Carrier radiation and ringdown dynamics & 215\\[9pt]
\textbf{XII} & Cosmological fate and the completed boundary & 224\\[9pt]
\textbf{XIII} & Home: information on an endless chain of two-dimensional closed spheres & 229\\[9pt]
\multicolumn{3}{l}{\textit{Appendix}}\\[6pt]\textbf{A} & Evidence, closures and remaining physical assignments & 235\\[9pt]
\bottomrule\end{tabularx}
\vspace{0.2in}
\noindent There is one volume-wide contents list and one continuous page sequence. Chapter and section bookmarks provide navigation within the PDF. Equations, figures, tables and references are numbered locally within their chapter; citations should therefore name the chapter as well as the local identifier.
\clearpage
\section*{How to read this edition}
\subsection*{A physical construction with its assumptions visible}
The volume moves from finite source histories and retained accounts to spatial response, clocks and signals, strong fields, quantum records, boundary thermodynamics, cosmic inventory, thermal history, halos, radiation and possible cosmic futures. The same word can describe different mathematical objects in these settings. A count is not automatically a mass; a propagation account is not automatically a density; a finite source calculation is not automatically a long-time cosmological law. Each join must identify the receiving law, its units, and the observable it supplies.

The expanded chapters separate exact identities, theorems under declared assumptions, results reported by retained computational records, and physical hypotheses awaiting selection or calibration. The new worked examples are analytic comparisons and finite arithmetic checks. They do not constitute new native GEN6 campaigns, new observational fits, or experimental validation. Where the submitted closing review claimed more than its evidence established, Appendix A records the narrower conclusion and the remaining interface.

\subsection*{Continuity with the submitted manuscript}
The latest submitted content, rather than an older shorter chapter, is the preservation baseline. Chapters I--III match their earlier formatting sources exactly. The integrated source for Chapter IV, Chapter V and Chapter VI was recovered from the archived research revision \texttt{6890eae4cb5f33c766c86ce10f43b5165a4c4157}. The original vector content of Chapters I--V is retained, with harmonized running headers and continuous folios. Chapter VI retains its complete source-derived body and proofs while changing its layout. Historical dates, receipt identifiers and claims of historical verification keep their original scope.

Some retained passages describe earlier development stages. Read them together with their later continuation and the current evidence appendix. A historical section's list of open tasks is not necessarily the final status of the chapter. Conversely, a conditional closure is not promoted to an unconditional physical law by inclusion in this volume. References to old packet pages identify those archived packets unless an explicit chapter and current volume page is given.

\subsection*{Sources, code and verification}
The companion source package contains the editable chapter sources, figures, manuscript build instructions, source-recovery provenance, new arithmetic checks and assembly verification. The evidence appendix inventories the earlier review's identifiers so that revised conclusions remain traceable. Historical native evidence and observational results are reported from their sources; they have not been represented as rerun for this edition.

The opening chapter's public data repository is \url{https://github.com/iwtbotiwtwot/sam_volume_I_i}; Chapter II's is \url{https://github.com/iwtbotiwtwot/sam_volume_I_ii}. Other chapter-specific source records and references remain with their respective chapters. The archived research revision used for source recovery is identified in the companion provenance manifest. Repository accessibility and the presence of a receipt do not replace checking its mathematical or empirical scope.

\subsection*{What the expansion is intended to accomplish}
The later chapters now explain the intermediate calculations that short drafts left implicit: thermodynamic balances and accessible information; the dimensionful inventory cascade; recombination and visibility; potential, force and halo inversion; radiative flux and spectral boundary conditions; and conservation, causal horizons and future evolution. Solved examples expose which conclusions follow from the stated equations and which require an additional physical law. This is a fuller working manuscript for scientific review, with no imposed chapter page limit.
\end{document}
